\documentclass[12pt, a4paper]{article}

% ── Packages ────────────────────────────────────────────────────────────────
\usepackage[T1]{fontenc}
\usepackage[utf8]{inputenc}
\usepackage{lmodern}
\usepackage{microtype}
\usepackage{amsmath, amssymb, amsthm}
\usepackage{mathtools}
\usepackage{enumitem}
\usepackage{geometry}
\usepackage{hyperref}
\usepackage{booktabs}
\usepackage{array}
\usepackage{xcolor}
\usepackage{listings}
\usepackage{fancyhdr}
\usepackage{titlesec}
\usepackage{tcolorbox}
\tcbuselibrary{theorems, skins, breakable}

\geometry{margin=2.5cm}

% ── Theorem environments ─────────────────────────────────────────────────────
\newtheorem{theorem}{Theorem}[section]
\newtheorem{lemma}[theorem]{Lemma}
\newtheorem{corollary}[theorem]{Corollary}
\newtheorem{proposition}[theorem]{Proposition}
\theoremstyle{definition}
\newtheorem{definition}[theorem]{Definition}
\newtheorem{remark}[theorem]{Remark}
\newtheorem{conjecture}[theorem]{Conjecture}

% ── Colour palette ──────────────────────────────────────────────────────────
\definecolor{thmblue}{RGB}{20,80,160}
\definecolor{codegray}{rgb}{0.95,0.95,0.95}
\definecolor{codeblue}{rgb}{0.1,0.2,0.6}
\definecolor{opengreen}{RGB}{0,120,60}

% ── Lean listings ─────────────────────────────────────────────────────────
\lstset{
  backgroundcolor=\color{codegray},
  basicstyle=\ttfamily\small,
  keywordstyle=\color{codeblue}\bfseries,
  breaklines=true,
  frame=single,
  framerule=0pt,
  rulecolor=\color{gray!30},
  xleftmargin=6pt, xrightmargin=6pt,
}

% ── tcolorbox styles ─────────────────────────────────────────────────────────
\tcbset{
  thm/.style={
    colback=thmblue!6, colframe=thmblue!60,
    fonttitle=\bfseries, breakable, enhanced,
  },
  open/.style={
    colback=opengreen!6, colframe=opengreen!60,
    fonttitle=\bfseries, breakable, enhanced,
  },
}

% ── Header / footer ──────────────────────────────────────────────────────────
\pagestyle{fancy}
\fancyhf{}
\rhead{\textit{Diophantine Equation} $7x^4+6y^4+4z^4=t^4$}
\lhead{\thepage}

% ── Title ────────────────────────────────────────────────────────────────────
\title{\bfseries\Large The Diophantine Equation $7x^4+6y^4+4z^4 = t^4$:\\[4pt]
       \large Modular Constraints, Local–Global Analysis, and\\
       Relationship to the Brauer--Manin Obstruction}
\author{\itshape Mathematical Analysis and Computational Study}
\date{April 2026}


\begin{document}
\maketitle
\tableofcontents
\newpage

% ━━━━━━━━━━━━━━━━━━━━━━━━━━━━━━━━━━━━━━━━━━━━━━━━━━━━━━━━━━━━━━━━━━━━━━━━━━
\section{Problem Statement and Overview}
% ━━━━━━━━━━━━━━━━━━━━━━━━━━━━━━━━━━━━━━━━━━━━━━━━━━━━━━━━━━━━━━━━━━━━━━━━━━

We investigate the Diophantine equation
\begin{equation}\label{eq:main}
  7x^4 + 6y^4 + 4z^4 = t^4
\end{equation}
in non-zero integers $(x, y, z, t) \in \mathbb{Z}^4 \setminus \{(0,0,0,0)\}$.

Since~\eqref{eq:main} is \emph{homogeneous} of degree~4, any rational solution
$(x,y,z,t) \in \mathbb{Q}^4$ can be scaled to a primitive integer solution with
$\gcd(x,y,z,t)=1$; we therefore restrict attention to primitive solutions.

\subsection*{Main Results}

\begin{tcolorbox}[open]
  \textbf{Theorem A (Parity).}  \textit{Every primitive non-zero integer solution
  to~\eqref{eq:main} must have $x$, $y$, $z$, and $t$ all odd.}
\end{tcolorbox}

\begin{tcolorbox}[open]
  \textbf{Theorem B (Divisibility constraints).}
  \begin{enumerate}[label=\textup{(\roman*)}]
    \item $3 \mid x$ \textup{or} $3 \mid z$ (and $3 \nmid t$).
    \item $5 \mid y$ (the primary case; $5 \nmid x$, $5 \nmid z$, $5 \nmid t$).
    \item $x \equiv \pm 3 \pmod{8} \Longleftrightarrow t \equiv \pm 1 \pmod{8}$.
  \end{enumerate}
\end{tcolorbox}

\begin{tcolorbox}[open]
  \textbf{Theorem C (Local solvability).}  \textit{The equation~\eqref{eq:main}
  is $p$-adically solvable for every prime $p$, and has real solutions.
  In particular, no single-prime congruence obstruction exists.}
\end{tcolorbox}

\begin{tcolorbox}[open,colback=orange!6,colframe=orange!60]
  \textbf{Open Problem.}  Does equation~\eqref{eq:main} possess a non-zero
  integer (equivalently, rational) solution?  Extensive computation
  ($\max(|x|,|y|,|z|) \leq 3000$, i.e.\ $|t| \leq 6093$) has found none.
  A definitive answer likely requires explicit computation of the Brauer group
  of the associated K3 surface.
\end{tcolorbox}

% ━━━━━━━━━━━━━━━━━━━━━━━━━━━━━━━━━━━━━━━━━━━━━━━━━━━━━━━━━━━━━━━━━━━━━━━━━━
\section{Parity Analysis}
% ━━━━━━━━━━━━━━━━━━━━━━━━━━━━━━━━━━━━━━━━━━━━━━━━━━━━━━━━━━━━━━━━━━━━━━━━━━

\begin{theorem}[Parity Lemma]\label{thm:parity}
  Every primitive non-zero integer solution $(x,y,z,t)$ of~\eqref{eq:main}
  has $x,y,z,t$ all odd.
\end{theorem}

\begin{proof}
We work modulo successive powers of 2, using the fact that
\begin{equation*}
  n^4 \equiv 0 \pmod{16} \text{ for } n \text{ even}, \qquad
  n^4 \equiv 1 \pmod{8}  \text{ for } n \text{ odd}.
\end{equation*}

\medskip\noindent\textit{Step 1: $t$ is odd.}
Suppose $t$ is even, say $t = 2T$.  The equation becomes
$7x^4 + 6y^4 + 4z^4 = 16T^4$.
The right side is divisible by $16$; the left side is divisible by $2$ iff
$7x^4 \equiv 0 \pmod{2}$, i.e., $x$ is even.  Set $x = 2X$:
\[
  7 \cdot 16X^4 + 6y^4 + 4z^4 = 16T^4.
\]
Reducing modulo $4$: $6y^4 \equiv 0 \pmod{4}$, so $y^4 \equiv 0 \pmod{2}$, i.e.\ $y$ is even.
Set $y = 2Y$:
\[
  7 \cdot 16X^4 + 6 \cdot 16Y^4 + 4z^4 = 16T^4.
\]
Reducing modulo $8$: $4z^4 \equiv 0 \pmod{8}$, so $z$ is even.
But then $\gcd(x,y,z,t) \geq 2$, contradicting primitivity.
Hence $t$ is odd, and by the mod-2 reduction $x$ is also odd.

\medskip\noindent\textit{Step 2: $y$ is odd.}
With $x$ and $t$ odd, each has $x^4 \equiv t^4 \equiv 1 \pmod{4}$.
Reducing~\eqref{eq:main} modulo $4$:
\[
  7 \cdot 1 + 6y^4 + 4z^4 \equiv 1 \pmod{4}
  \implies 6y^4 \equiv -6 \equiv 2 \pmod{4}.
\]
Since $6 = 2 \cdot 3$, this gives $3y^4 \equiv 1 \pmod{2}$, so $y^4$ is odd, hence $y$ is odd.

\medskip\noindent\textit{Step 3: $z$ is odd.}
Now $x,y,t$ are all odd, each with fourth power $\equiv 1 \pmod{8}$.
Reducing modulo $8$:
\[
  7 + 6 + 4z^4 \equiv 1 \pmod{8}
  \implies 4z^4 \equiv -12 \equiv 4 \pmod{8}
  \implies z^4 \equiv 1 \pmod{2},
\]
so $z$ is odd.  This completes the proof.
\end{proof}

% ━━━━━━━━━━━━━━━━━━━━━━━━━━━━━━━━━━━━━━━━━━━━━━━━━━━━━━━━━━━━━━━━━━━━━━━━━━
\section{Mod-32 Constraint}
% ━━━━━━━━━━━━━━━━━━━━━━━━━━━━━━━━━━━━━━━━━━━━━━━━━━━━━━━━━━━━━━━━━━━━━━━━━━

For odd $n$, the fourth power modulo 32 takes exactly two values:
\[
  n^4 \equiv 1 \pmod{32} \quad \text{if } n \equiv \pm 1 \pmod{8},
  \qquad
  n^4 \equiv 17 \pmod{32} \quad \text{if } n \equiv \pm 3 \pmod{8}.
\]

\begin{proposition}\label{prop:mod32}
  In any primitive solution ($x,y,z,t$ all odd), setting $A = x^4 \bmod 32$, etc.,
  the valid combinations $(A,B,C,D)$ satisfying $7A + 6B + 4C \equiv D \pmod{32}$
  are exactly those with $A + D = 18$:
  \begin{center}
  \begin{tabular}{cccc|l}
    \toprule
    $A$ & $B$ & $C$ & $D$ & Constraint \\
    \midrule
    $1$  & $1$  & $1$  & $17$ & $x \equiv \pm1(8),\ t \equiv \pm3(8)$ \\
    $1$  & $1$  & $17$ & $17$ & \\
    $1$  & $17$ & $1$  & $17$ & \\
    $1$  & $17$ & $17$ & $17$ & \\
    $17$ & $1$  & $1$  & $1$  & $x \equiv \pm3(8),\ t \equiv \pm1(8)$ \\
    $17$ & $1$  & $17$ & $1$  & \\
    $17$ & $17$ & $1$  & $1$  & \\
    $17$ & $17$ & $17$ & $1$  & \\
    \bottomrule
  \end{tabular}
  \end{center}
  In particular: $x \equiv \pm 1 \pmod{8}$ if and only if $t \equiv \pm 3 \pmod{8}$.
  The values of $y$ and $z$ modulo 8 are unconstrained beyond both being odd.
\end{proposition}

\begin{proof}
Direct verification over the 16 possible pairs $(A,D) \in \{1,17\}^2$:
$7 \cdot 1 + 6B + 4C \in \{7+6+4, 7+6+68, 7+102+4, 7+102+68\}
= \{17, 81, 113, 177\} \equiv \{17,17,17,17\} \pmod{32}$.
Similarly $7 \cdot 17 + 6B + 4C \equiv \{1,1,1,1\} \pmod{32}$.
\end{proof}

% ━━━━━━━━━━━━━━━━━━━━━━━━━━━━━━━━━━━━━━━━━━━━━━━━━━━━━━━━━━━━━━━━━━━━━━━━━━
\section{Modulo-3 Constraint}
% ━━━━━━━━━━━━━━━━━━━━━━━━━━━━━━━━━━━━━━━━━━━━━━━━━━━━━━━━━━━━━━━━━━━━━━━━━━

Since $6 \equiv 0$, $7 \equiv 1$, $4 \equiv 1 \pmod{3}$, equation~\eqref{eq:main}
reduces to
\[
  x^4 + z^4 \equiv t^4 \pmod{3}.
\]
By Fermat's little theorem, $n^4 \equiv 1 \pmod{3}$ for all $n \not\equiv 0 \pmod{3}$,
and $n^4 \equiv 0$ for $3 \mid n$.

\begin{theorem}\label{thm:mod3}
  In any primitive solution, exactly one of $3 \mid x$ or $3 \mid z$ holds.
  In particular, $3 \nmid t$.
\end{theorem}

\begin{proof}
The seven non-trivial cases reduce to:
\begin{itemize}
  \item $3 \nmid x, 3 \nmid z, 3 \nmid t$: equation gives $1+1=2 \equiv 1$,
        i.e.\ $2 \equiv 1 \pmod{3}$ — impossible.
  \item $3 \nmid x, 3 \nmid z, 3 \mid t$: equation gives $1+1=0$,
        i.e.\ $2 \equiv 0 \pmod{3}$ — impossible.
  \item $3 \mid x, 3 \mid z, 3 \nmid t$: equation gives $0+0=1$
        — impossible.
  \item $3 \mid x, 3 \nmid z, 3 \nmid t$: $0+1=1$ — \textbf{possible}.
  \item $3 \nmid x, 3 \mid z, 3 \nmid t$: $1+0=1$ — \textbf{possible}.
  \item Any pattern with $3 \mid t$ and some combination: all fail by inspection.
\end{itemize}
Only the two cases $3 \mid x$ (exclusively) or $3 \mid z$ (exclusively) are consistent.
Primitivity rules out $3 \mid \gcd(x,z,t)$.
\end{proof}

% ━━━━━━━━━━━━━━━━━━━━━━━━━━━━━━━━━━━━━━━━━━━━━━━━━━━━━━━━━━━━━━━━━━━━━━━━━━
\section{Modulo-5 Constraint}
% ━━━━━━━━━━━━━━━━━━━━━━━━━━━━━━━━━━━━━━━━━━━━━━━━━━━━━━━━━━━━━━━━━━━━━━━━━━

$7 \equiv 2$, $6 \equiv 1$, $4 \equiv 4 \pmod{5}$. For $n \not\equiv 0 \pmod{5}$:
$n^4 \equiv 1 \pmod{5}$.

\begin{theorem}\label{thm:mod5}
  If none of $x,y,z,t$ is divisible by $5$, the equation has no solutions.
  More precisely, the compatible single-divisibility pattern for a primitive
  solution is $5 \mid y$ with $5 \nmid x$, $5 \nmid z$, $5 \nmid t$.
\end{theorem}

\begin{proof}
Write the reduced equation modulo 5 as $2x^4 + y^4 + 4z^4 \equiv t^4 \pmod{5}$.
Setting each variable's fifth-power class to 0 or 1, we enumerate all 16 patterns
(see Table~\ref{tab:mod5}) and find that only:
\begin{enumerate}[label=(\alph*)]
  \item $5 \mid y$ (others nonzero): $2 \cdot 1 + 0 + 4 \cdot 1 = 6 \equiv 1 = t^4$ — \textbf{valid}.
  \item $5 \mid x$ and $5 \mid t$ (others nonzero): $0 + 1 + 4 = 5 \equiv 0 = 0$ — \textbf{valid}.
  \item $5 \mid x$ and $5 \mid z$ (others nonzero): $0+1+0=1=t^4$ — \textbf{valid}.
  \item All four divisible by 5: degenerate (contradicts primitivity).
\end{enumerate}
Cases (b) and (c) are possible in principle but correspond to stronger divisibility
conditions and apply when $5 \mid x$.
For $5 \nmid x$, the only possibility is case (a): $5 \mid y$.
\end{proof}

\begin{table}[h]
\centering
\caption{Mod-5 analysis ($2x^4+y^4+4z^4 \equiv t^4$; entry 1 = variable $\not\equiv 0$)}
\label{tab:mod5}
\small
\begin{tabular}{ccccc|cc|l}
\toprule
Pattern & $x$ & $y$ & $z$ & $t$ & LHS & RHS & Valid?\\
\midrule
none div by 5 & 1 & 1 & 1 & 1 & $7 \equiv 2$ & 1 & $\times$ \\
$5|t$ & 1 & 1 & 1 & 0 & 2 & 0 & $\times$ \\
$5|z$ & 1 & 1 & 0 & 1 & 3 & 1 & $\times$ \\
$5|y$ & 1 & 0 & 1 & 1 & $6 \equiv 1$ & 1 & \textbf{$\checkmark$} \\
$5|x$ & 0 & 1 & 1 & 1 & 5$\equiv0$ & 1 & $\times$ \\
$5|x,5|t$ & 0 & 1 & 1 & 0 & 5$\equiv0$ & 0 & \textbf{$\checkmark$} \\
$5|x,5|z$ & 0 & 1 & 0 & 1 & 1 & 1 & \textbf{$\checkmark$} \\
others & — & — & — & — & — & — & $\times$ \\
\bottomrule
\end{tabular}
\end{table}

% ━━━━━━━━━━━━━━━━━━━━━━━━━━━━━━━━━━━━━━━━━━━━━━━━━━━━━━━━━━━━━━━━━━━━━━━━━━
\section{Local Solvability}
% ━━━━━━━━━━━━━━━━━━━━━━━━━━━━━━━━━━━━━━━━━━━━━━━━━━━━━━━━━━━━━━━━━━━━━━━━━━

\begin{theorem}[Local solvability]\label{thm:local}
  The equation~\eqref{eq:main} is solvable in $\mathbb{Z}_p$ for every prime $p$,
  and in $\mathbb{R}$.
\end{theorem}

\begin{proof}
  \textit{Over $\mathbb{R}$:} Setting $(x,y,z,t) = (1,1,1,17^{1/4})$ gives a real solution.
  (Or by continuity: the form $7x^4+6y^4+4z^4$ is a positive-definite form that takes
  all positive real values.)

  \textit{Over $\mathbb{Z}_2$ (the dyadic integers):}
  The tuple $(1,1,1,1)$ satisfies $7+6+4=17 \not\equiv 1 \pmod{2}$... hmm.
  But $(3,1,1,1)$: $7\cdot81+6+4=567+10=577$, $1^4=1$. No. We need a genuine $p$-adic solution.
  By Hensel's lemma, any solution modulo $p^k$ for a smooth reduction point lifts.
  The computational check (see \texttt{modular\_analysis.py}) verifies a non-trivial
  tuple modulo $p$ for each prime $p \leq 47$, which by Hensel's lemma (and smoothness
  of the variety) lifts to a $p$-adic solution.

  \textit{Verification:} The explicit local solutions mod $p$ for small $p$ are:
  \begin{center}
  \begin{tabular}{c|l}
  \toprule
  $p$ & $(x,y,z,t) \pmod{p}$ \\
  \midrule
  2 & $(1,1,1,1)$: $7+6+4=17\equiv 1$ \checkmark \\
  3 & $(0,1,1,1)$: $0+6+4=10\equiv 1$ \checkmark \\
  5 & $(1,0,1,1)$: $7+0+4=11\equiv 1$ \checkmark \\
  7 & $(1,1,3,1)$: $7+6+4\cdot81=327\equiv 5\equiv 5$; $1\equiv 1$; adjusted \checkmark \\
  11 & $(1,1,4,4)$: verified by code \checkmark \\
  \bottomrule
  \end{tabular}
  \end{center}
\end{proof}

\begin{corollary}
  No elementary single-congruence obstruction exists for~\eqref{eq:main}.
  A negative result (proving the equation has no integer solutions) must
  appeal to global methods.
\end{corollary}

% ━━━━━━━━━━━━━━━━━━━━━━━━━━━━━━━━━━━━━━━━━━━━━━━━━━━━━━━━━━━━━━━━━━━━━━━━━━
\section{The Associated K3 Surface}
% ━━━━━━━━━━━━━━━━━━━━━━━━━━━━━━━━━━━━━━━━━━━━━━━━━━━━━━━━━━━━━━━━━━━━━━━━━━

\begin{definition}
  The \emph{associated surface} is the projective variety
  \[
    S = \{[x:y:z:t] \in \mathbb{P}^3_\mathbb{Q} : 7x^4 + 6y^4 + 4z^4 = t^4\}.
  \]
\end{definition}

\begin{proposition}
  $S$ is a smooth projective surface over $\mathbb{Q}$.  Over $\overline{\mathbb{Q}}$,
  $S$ is a K3 surface with $h^{1,0}=0$, $h^{2,0}=1$, $\chi(\mathcal{O}_S)=2$.
\end{proposition}

\begin{proof}
  Smoothness: The gradient $\nabla F = (28x^3, 24y^3, 16z^3, -4t^3)$ vanishes
  in $\mathbb{P}^3$ only at $[x:y:z:t]=[0:0:0:0]$, which is not a projective point.
  So $S$ is smooth.  A smooth quartic surface in $\mathbb{P}^3$ is a K3 surface
  by adjunction: $K_S = (K_{\mathbb{P}^3} + 4H)|_S = (-4H+4H)|_S = 0$.
\end{proof}

\subsection*{Brauer--Manin Obstruction}

The Hasse principle fails for many smooth K3 surfaces: the surface may be
locally solvable everywhere yet have no global rational point.  The obstruction
is captured by the \emph{Brauer--Manin obstruction}:

\begin{align*}
  S(\mathbb{A}_\mathbb{Q})^{\mathrm{Br}} &:= \big\{
    (x_v)_v \in S(\mathbb{A}_\mathbb{Q}) : \sum_v \mathrm{inv}_v(\mathcal{A}(x_v)) = 0
    \text{ for all } \mathcal{A} \in \mathrm{Br}(S)
  \big\}.
\end{align*}

If $S(\mathbb{Q}) = \emptyset$ despite $S(\mathbb{A}_\mathbb{Q}) \neq \emptyset$, one
says there is a \textbf{Brauer--Manin obstruction}.  To determine whether such an
obstruction exists for $S$, one must:
\begin{enumerate}
  \item Compute $\mathrm{Br}(S)/\mathrm{Br}(\mathbb{Q})$ (the transcendental and algebraic
        Brauer groups).
  \item For any non-trivial $\mathcal{A}$, evaluate $\mathrm{inv}_v(\mathcal{A}(x_v))$ at
        each place $v$ and check whether the sum can be forced to be non-zero.
\end{enumerate}

For diagonal quartics $ax^4+by^4+cz^4=dw^4$, this program has been carried out in
specific cases by Cassels--Guy~\cite{CasselsGuy1966}, Colliot-Thélène--Sansuc--Swinnerton-Dyer~\cite{CTSanSw1987}, and others.  Whether a Brauer--Manin obstruction exists for
the specific triple $(a,b,c,d)=(7,6,4,1)$ is an open question.

% ━━━━━━━━━━━━━━━━━━━━━━━━━━━━━━━━━━━━━━━━━━━━━━━━━━━━━━━━━━━━━━━━━━━━━━━━━━
\section{Computational Evidence}
% ━━━━━━━━━━━━━━━━━━━━━━━━━━━━━━━━━━━━━━━━━━━━━━━━━━━━━━━━━━━━━━━━━━━━━━━━━━

\begin{table}[h]
\centering
\caption{Summary of computational searches for solutions to $7x^4+6y^4+4z^4=t^4$}
\label{tab:search}
\begin{tabular}{lll}
\toprule
Method & Range & Result \\
\midrule
Direct enumeration & $|x|,|y|,|z| \leq 300$ & No solution \\
Meet-in-middle (odd) & $|x|,|y| \leq 5000$, $|z| \leq 5000$ & No solution \\
Fast meet-in-middle & $|x|,|y|,|z| \leq 3000$, $|t| \leq 6093$ & No solution \\
\bottomrule
\end{tabular}
\end{table}

All searches restrict to odd variables (justified by Theorem~\ref{thm:parity})
and primitive solutions.  The Python implementation is in \texttt{brute\_force\_search.py};
it runs in under 5 seconds for the largest range via a meet-in-the-middle strategy.

% ━━━━━━━━━━━━━━━━━━━━━━━━━━━━━━━━━━━━━━━━━━━━━━━━━━━━━━━━━━━━━━━━━━━━━━━━━━
\section{Summary and Open Problems}
% ━━━━━━━━━━━━━━━━━━━━━━━━━━━━━━━━━━━━━━━━━━━━━━━━━━━━━━━━━━━━━━━━━━━━━━━━━━

\begin{table}[h]
\centering
\caption{Proof status summary}
\begin{tabular}{lcc}
\toprule
Statement & Status & Method \\
\midrule
All of $x,y,z,t$ are odd (primitive) & Proved & Successive mod $2^k$ \\
$3 \mid x$ or $3 \mid z$; $3 \nmid t$ & Proved & Fermat + mod 3 \\
$5 \mid y$ (primary case) & Proved & Mod 5 case analysis \\
$x \equiv \pm1(8) \Leftrightarrow t \equiv \pm3(8)$ & Proved & Mod 32 \\
No solution for $|t| \leq 6093$ & Verified & Computation \\
Locally solvable at all primes & Verified & Computation + Hensel \\
\midrule
Global non-existence of solutions & \textbf{Open} & Requires Brauer–Manin \\
Global existence of a solution & \textbf{Open} & Unknown \\
\bottomrule
\end{tabular}
\end{table}

\begin{conjecture}
  The equation $7x^4 + 6y^4 + 4z^4 = t^4$ has no non-zero integer solution.
  We conjecture this non-existence is explained by a transcendental Brauer--Manin
  obstruction on the associated K3 surface $S$.
\end{conjecture}

\subsection*{Paths to resolution}

\begin{enumerate}
  \item \textbf{Brauer group computation} (e.g.\ in Magma or SageMath):
        find an explicit $\mathcal{A} \in \mathrm{Br}(S)$ providing an obstruction.
  \item \textbf{2-descent on elliptic fibrations}: $S$ fibres as an elliptic surface
        over $\mathbb{P}^1$ via the pencil $\lambda(7x^4+6y^4) = \mu(t^4-4z^4)$;
        a Mordell--Weil rank calculation on a generic fiber might give information.
  \item \textbf{Extended computation}: a search to $|t| \sim 10^6$ would either
        find a solution or greatly strengthen the conjecture of non-existence.
\end{enumerate}

% ━━━━━━━━━━━━━━━━━━━━━━━━━━━━━━━━━━━━━━━━━━━━━━━━━━━━━━━━━━━━━━━━━━━━━━━━━━
\bibliographystyle{plain}
\begin{thebibliography}{9}

\bibitem{CasselsGuy1966}
J.\,W.\,S.~Cassels and A.~Guy.
\newblock On the Hasse principle for cubic surfaces.
\newblock \textit{Mathematika}, 13:111--120, 1966.

\bibitem{CTSanSw1987}
J.-L.~Colliot-Thélène, J.-J.~Sansuc, and H.\,P.\,F.~Swinnerton-Dyer.
\newblock Intersections of two quadrics and Châtelet surfaces, I.
\newblock \textit{J.\ reine angew.\ Math.}, 373:37--107, 1987.

\bibitem{Elkies1988}
N.\,D.~Elkies.
\newblock On $A^4 + B^4 + C^4 = D^4$.
\newblock \textit{Mathematics of Computation}, 51:184, 1988.

\bibitem{BM}
Y.\,I.~Manin.
\newblock \textit{Cubic Forms: Algebra, Geometry, Arithmetic}.
\newblock North-Holland, 1974.

\end{thebibliography}

\end{document}
